\input{../tex-inputs/latex-leseansicht-vorspann}

\section[Gustav Schwarzkopf an Arthur Schnitzler, 19. 5. 1909]{L04319 Gustav Schwarzkopf an Arthur Schnitzler, 19. 5. 1909}
\nopagebreak\mylabel{L04319v}
\rehead{ }\normalsize\beginnumbering\briefempfaengerindex{Schnitzler, Arthur@\textsc{Schnitzler, Arthur}!zzzSchwarzkopf, Gustav@\emph{von Gustav Schwarzkopf}!1909-05-191@{19. 5. 1909}|(be}
\toendnotes[C]{\smallbreak\pagebreak[2]}
\correspDesc{Versand  durch Gustav Schwarzkopf am 19. 5. 1909 in Wien
\newline{}Übermittlung  am 19. 5. 1909 in Wien
\newline{}Erhalt  durch Arthur Schnitzler im Zeitraum [19. 5. 1909 – 22. 5. 1909?] in Wien}\toendnotes[C]{\smallbreak}
\Standort{CUL, Schnitzler, B 96a.}
\physDesc{Postkarte, 555 Zeichen
\newline{}Handschrift: schwarze Tinte, deutsche Kurrent
\newline{}Versand: Stempel: »\nobreak{}\oindex{I., Innere Stadt@\textbf{I., Innere Stadt}, \emph{Verwaltungsgebiet}|pwk}1/1 Wien 1, 19. V. 09, 12\nobreak{}«.  
\newline{}Schnitzler: mit Bleistift Vermerk: »\textsc{Schwarzkopf}« }\toendnotes[C]{\smallbreak}\pstart{}{\pb}I. Tiefer Graben 17\oindex{Wien@\textbf{Wien}!I., Innere Stadt@\textbf{I., Innere Stadt}!Tiefer Graben 17@\textbf{Tiefer Graben 17}, \emph{Wohngebäude}|pw}.\pend{}\pstart{}II. St.\hspace*{1.5em}Th. 6.\pend{}{\bigskip}\pstart{}Herrn\pend{}\pstart{}\textsc{D\textsuperscript{r} Arthur
                  Schnitzler}\pend{}\pstart{}\textsc{Wien\oindex{Wien@\textbf{Wien}, \emph{Verwaltungsgebiet}|pw}}\pend{}\pstart{}XVIII. Spöttelgaſſe 7\oindex{Wien@\textbf{Wien}!XVIII., Währing@\textbf{XVIII., Währing}!Edmund-Weiß-Gasse@\textbf{Edmund-Weiß-Gasse}, \emph{Straße}|pw}.\pend{}{\bigskip}\vspace{1em}
\pstart
           \raggedleft{}{\pb}19. V. 09\pend
           \vspace{0.5em}
\pstart
           Lieber Arthur, es iſt wirklich{ }ſehr lieb, daß Sie{ }ſich in dieſer
               Sache{ }ſo bemühen. Da Sie das kleine Mädchen\pwindex{?? [Mädchen, dem Schnitzler im Mai 1909 hilft] @\textsc{?? [Mädchen, dem Schnitzler im Mai 1909 hilft]}|pwv} nicht ke{\geminationn}en, muß ich Ihre Bemühung als eine \uline{mir} erwieſene Gefälligkeit anſehen und danke Ihnen beſtens. Vor einem
               begeiſterten Dankſchreiben von Tante\pwindex{?? [Tante von Mädchen, dem Schnitzler hilft, 1909] @\textsc{?? [Tante von Mädchen, dem Schnitzler hilft, 1909]}|pwv}
               und Nichte\pwindex{?? [Mädchen, dem Schnitzler im Mai 1909 hilft] @\textsc{?? [Mädchen, dem Schnitzler im Mai 1909 hilft]}|pwv} habe ich Sie bis jetzt
               geſchützt, aber auf die Dauer wird es{ }ſich wol nicht verhindern laſſen. Einſtweilen
               iſt man von Ihrer »großen Güte« einfach entzückt.\pend
           
\pstart
           Ihre Frau\pwindex{Schnitzler, Olga 17.\,1.\,1882 Wien – 13.\,1.\,1970 Lugano@\textsc{Schnitzler, Olga} (17.\,1.\,1882 Wien – 13.\,1.\,1970 Lugano), \emph{Schauspielerin, Sängerin}|pwv} und Sie
               herzlichſt grüßend{\\[\baselineskip]}Ihr \spacefill\mbox{Gustav.}\pend
           \leftskip=0em{}\selectlanguage{ngerman}\endnumbering\briefempfaengerindex{Schnitzler, Arthur@\textsc{Schnitzler, Arthur}!zzzSchwarzkopf, Gustav@\emph{von Gustav Schwarzkopf}!1909-05-191@{19. 5. 1909}|)be}\mylabel{L04319h}
\begin{anhang}
\end{anhang}\newcommand{\dateiname}{L04319}\newcommand{\titel}{Gustav Schwarzkopf an Arthur Schnitzler, 19. 5. 1909}\newcommand{\editorInnen}{Herausgegeben von Jahnke, SelmaMüller, Martin Anton}\input{../tex-inputs/latex-leseansicht-abspann}
